% Additional paragraph for Section 8.3
% Addresses how k=4 closure works for composite objects (e.g., asteroid with internal structure)
% Insert after equation (\ref{eq:poisson-numeric}) and before "This is what self-similarity means..."

An important mechanical question arises for composite objects: when an asteroid contains $10^{33}$ atoms, each with its own field structure, how does the unified gravitational envelope maintain $k=4$ closure? The field cannot be a simple average—this would wash out the topological structure—nor can it be purely isomorphic to the physical mass distribution (which may be highly irregular). The resolution is that field closure operates at \emph{radial domain boundaries}: the transition surfaces where field lines connect nested field envelopes at different scales. At each such boundary, the field must satisfy $\det(e) \neq 0$ with $k=4$ tetrahedral basis structure, regardless of the physical asymmetry of the object. This is the only way to achieve continuous field closure across scale transitions. The asteroid's envelope field is $k=4$ compliant at its outer boundary (the gravitational radius where the field transitions to the external vacuum), and internally, the field maintains $k=4$ structure at the radial boundaries separating atomic-scale envelopes from the macroscopic envelope. These boundaries are not physical surfaces but mathematical transition regions where the field's characteristic scale changes. The $k=4$ topology ensures that field closure holds at every such transition, allowing nested hierarchies without requiring the internal structure to be spherically symmetric or averaged. The physical asteroid may be elongated and irregular, but the field's topological structure—encoded in the tetrad bases $e_\mu^a$ satisfying $\det(e) \neq 0$—is constrained to have $k=4$ faces at every point in spacetime, including at all radial domain boundaries where scale transitions occur.
